% Scope: Derive coil encoding, SENSE, GRAPPA and noise amplification; separate acceleration geometry from reconstruction assumptions.
% Status: architecture only; no chapter prose or completed problems yet.
% Prerequisites: Chapters 1, 5 and 6; receive hardware in Chapter 17.
% Planned worked examples: Two-coil SENSE inversion and g-factor SNR penalty.
% Planned figures: Aliased voxels, sensitivity vectors and GRAPPA neighborhoods.
% Planned challenge problems: Find an encoding null space; compare calibration and noise penalties.
\chapter{Parallel Imaging}
\label{ch:14}

\section{Receive arrays as encoding systems}
\label{sec:14:receive-arrays-as-encoding-systems}
\subsection{Sensitivity fields and folded voxels}
\subsection{Noise covariance and prewhitening}
\subsection{Optimal coil combination and normalization}

\section{SENSE reconstruction}
\label{sec:14:sense-reconstruction}
\subsection{Local unfolding equations}
\subsection{Weighted least squares and conditioning}
\subsection{Noise propagation and g-factor}

\section{GRAPPA reconstruction}
\label{sec:14:grappa-reconstruction}
\subsection{Autocalibration and neighborhood kernels}
\subsection{Calibration fits and synthesized samples}
\subsection{Noise correlations and calibration failure}

\section{Simultaneous multislice reconstruction}
\label{sec:14:simultaneous-multislice-reconstruction}
\subsection{Slice separation and coil geometry}
\subsection{CAIPIRINHA shifts and controlled aliasing}
\subsection{Slice leakage and reconstruction trade-offs}

\section{Practical acceleration limits}
\label{sec:14:practical-acceleration-limits}
\subsection{Calibration overhead and net acceleration}
\subsection{Motion and sensitivity mismatch}
\subsection{Regularization, residual aliasing and SNR efficiency}

\section{Worked examples}
\label{sec:14:worked-examples}
% Planned calculations: Two-coil SENSE inversion and g-factor SNR penalty.
\subsection{Derivation and quantitative calculation}
\subsection{Assumptions, limiting cases and scanner interpretation}

\section{Challenge problems}
\label{sec:14:challenge-problems}
% Problem design: Find an encoding null space; compare calibration and noise penalties.
% Use challengeproblem and stable labels prob:14:<topic>.
% Complete solutions belong in Appendix E, section for this chapter.
% Use \begin{solution}{prob:14:<topic>} ... \end{solution} there.
\subsection{Derivation and quantitative design}
\subsection{Model criticism and integrated reasoning}
